\section{Datenbeschaffung und Datenaufbereitung}

\subsection{Zielsetzung}
Ziel dieses Projektteils ist die strukturierte Beschaffung und Aufbereitung von Daten, um die Datenbank sowie 
das darauf aufbauende System zu testen. Der Fokus liegt dabei insbesondere auf der Transformation von Rohdaten in 
ein konsistentes und qualitätsgesichertes Format.

\subsection{Datenquelle}
Als primäre Datenquelle dienen öffentlich zugängliche Datensätze der Plattform Kaggle \cite{Kaggle_Synthetic_Person}. 
Die Plattform bietet eine Vielzahl an Datensammlungen aus unterschiedlichen Anwendungsbereichen und ist daher gut 
für datengetriebene Projekte geeignet.
Die Datensätze liegen in der Regel als JSON- oder CSV-Dateien vor und enthalten bereits strukturierte Informationen. 
Dennoch existiert kein Datensatz, der direkt auf unser Datenbankschema passt. Zudem handelt es sich häufig um Rohdaten, 
die Unstimmigkeiten, fehlende Werte oder uneinheitliche Formate enthalten.
Ergänzend wurden weitere Datensätze von Mendeley \cite{Mendeley_Dataset} und Hugging Face \cite{HuggingFace_Dataset} 
verwendet, um die Datenbasis zu erweitern.

\subsection{Datenaufbereitung}
Ein zentraler Schritt nach der Erstellung der Datenbank ist die Befüllung mit geeigneten Daten. Da 
keine Datensätze existieren, die direkt dem Datenbankschema entsprechen, müssen die vorhandenen 
Rohdaten entsprechend angepasst werden.

Der erste Schritt besteht in der Bereinigung der Daten. Dabei werden insbesondere fehlende Werte, 
inkonsistente Schreibweisen sowie weitere Unstimmigkeiten behandelt. Dies umfasst unter anderem das 
Vereinheitlichen von Duplikaten sowie das Entfernen unvollständiger Datensätze.

Im Anschluss erfolgt die Transformation der bereinigten Daten, um sie an die Anforderungen des Datenbankschemas 
anzupassen. Hierbei werden Datentypen entsprechend konvertiert und Attribute gegebenenfalls aufgeteilt oder 
zusammengeführt. Ziel dieses Schrittes ist es, eine einheitliche und logisch strukturierte Datenbasis zu schaffen.

Vor der weiteren Verarbeitung werden die Daten abschließend validiert, um eine hohe Datenqualität sicherzustellen 
und Probleme beim Einfügen in die Datenbank sowie bei späteren Analysen zu vermeiden.

\subsection{Vorbereitung für die Speicherung}
Nach der Aufbereitung werden die Daten in eine strukturierte Form überführt, die eine effiziente und 
möglichst einfache Integration in die Datenbank ermöglicht. Dabei wird bereits darauf geachtet, dass die 
Daten optimal an das zugrunde liegende Datenbankschema angepasst sind, um den Einfügeprozess zu vereinfachen.

\subsection{Datenimport}
Für den Import der Daten in die Datenbank wurde ein automatisierter Ansatz mittels Python-Skripten gewählt. 
Python eignet sich aufgrund seiner umfangreichen Bibliotheken im Bereich der Datenverarbeitung besonders gut 
für diesen Anwendungsfall.

Der Importprozess gliedert sich in folgende Schritte:
\begin{itemize}
\item Einlesen der Quelldateien
\item Erste Strukturprüfung der Daten
\item Transformation in geeignete Datenstrukturen
\end{itemize}

Dabei wird bewusst eine Trennung zwischen Datenaufbereitung und Datenimport eingehalten, um eine klare 
Modularisierung zu gewährleisten.

\subsection{Technologischer Ansatz}
Für die Umsetzung der Datenverarbeitung und -aufbereitung wurde Python verwendet. Ausschlaggebend hierfür 
sind die gute Unterstützung im Bereich der Datenverarbeitung, die große Auswahl an verfügbaren Bibliotheken 
sowie die einfache Erweiterbarkeit.

Konkrete Entscheidungen bezüglich der eingesetzten Datenbanktechnologien und weiterer Tools werden in einem 
separaten Abschnitt erläutert.

\subsection{Herausforderungen}
Im Rahmen der Datenaufbereitung traten mehrere Herausforderungen auf. Eine der größten bestand darin, 
geeignete und realistische Datensätze zu finden. Da der Fokus auf praxisnahen Daten lag, wurden insbesondere 
Datensätze mit Informationen zu realen Medikamenten oder Krankheiten benötigt.

Solche spezifischen Datensätze sind jedoch nur begrenzt verfügbar, wodurch die Suche zeitaufwendig war 
und teilweise Kompromisse bei der Auswahl getroffen werden mussten.
